% Part: first-order-logic
% Chapter: completeness
% Section: outline

\documentclass[../../../include/open-logic-section]{subfiles}

\begin{document}

\iftag{FOL}
      {\olfileid{fol}{com}{out}}
      {\olfileid{pl}{com}{out}}

\olsection{De route door het bewijs}

Het bewijs van de volledigheidsstelling heeft nogal wat stappen. Bij
een eerste lezing raak je daardoor gemakkelijk de draad kwijt. Daarom
zetten we eerst de hele route op een rij. We beginnen door anders naar
de vraag te kijken. De formulering ``als $\Gamma \Entails !A$, dan
$\Gamma \Proves !A$'' geeft nog weinig houvast. Om
$\Gamma \Proves !A$ te bewijzen, moeten we !!a{derivation} vinden. De
aanname $\Gamma \Entails !A$ vertelt niet meteen hoe we die afleiding
moeten bouwen. Bij sommige bewijssystemen kan dat rechtstreeks, maar
hier nemen we een andere route. Volledigheid zegt namelijk ook: ``als
$\Gamma$ consistent is, dan is zij vervulbaar.'' Nu kunnen we de
informatie in~$\Gamma$ als bouwmateriaal gebruiken. Samen met de
aanname dat zij consistent is, proberen we
\iftag{FOL}{!!a{structure}}{!!a{valuation}} te bouwen die elke
\iftag{FOL}{!!{sentence}}{!!{formula}} in~$\Gamma$ waar maakt. We weten
welk soort \iftag{FOL}{!!{structure}}{!!{valuation}} we nodig hebben:
precies een die doet wat $\Gamma$ beschrijft!{}

Bevat $\Gamma$ alleen \iftag{FOL}{atomaire
  !!{sentence}s}{!!{propositional variable}s}? Dan kunnen we er
gemakkelijk een model voor bouwen.\iftag{FOL}{ Stel dat alle atomaire
  !!{sentence}s de vorm $\Atom{P}{a_1,\dots,a_n}$ hebben. De $a_i$ zijn
  daarbij !!{constant}s.}{} We hoeven alleen het volgende te kiezen:
\iftag{FOL}{%
  !!a{domain}~$\Domain{M}$ en een betekenis voor~$P$ waarvoor
  $\Sat{M}{\Atom{P}{a_1,\ldots,a_n}}$. Daar is een eenvoudig recept
  voor. Neem $\Domain{M} = \Nat$ en $\Assign{\Obj c_i}{M} = i$. Staat
  $\Atom{P}{a_1,\ldots,a_n} \in \Gamma$, zet dan het tupel $\tuple{k_1,
    \dots, k_n}$ in $\Assign{P}{M}$. Hier is $k_i$ de index van het
  constante symbool~$a_i$ (dus $a_i \ident \Obj c_{k_i}$).}{%
  !!a{valuation}~$\pAssign{v}$ waarvoor $\pSat{v}{p}$ voor alle $p \in
  \Gamma$. Definieer $\pAssign{v}(p) = \True$ precies als
  $p \in \Gamma$.}

Neem nu aan dat $\Gamma$ een !!{formula}~$\lnot !B$ bevat en dat $!B$
atomair is. Kan onze bouw van
$\iftag{FOL}{\Struct{M}}{\pAssign{v}}$ dan verhinderen dat $\lnot !B$
waar wordt? Hier helpt de consistentie van~$\Gamma$. Als
$\lnot !B \in \Gamma$, dan $!B \notin \Gamma$; anders zou $\Gamma$
inconsistent zijn. En als $!B \notin \Gamma$, dan zorgt onze
constructie van~$\iftag{FOL}{\Struct{M}}{\pAssign{v}}$ ervoor dat
$\iftag{FOL}{\Sat/{M}}{\pSat/{v}}{!B}$. Dus
$\iftag{FOL}{\Sat{M}}{\pSat{v}}{\lnot !B}$. Dit geval werkt.

Maar wat doen we als $\Gamma$ samengestelde, niet-atomaire formules
bevat? Stel dat $!A \land !B$ erin staat. Om die formule waar te maken,
moeten we doen alsof zowel $!A$ als $!B$ in~$\Gamma$ staat. En als
$!A \lor !B \in \Gamma$, moeten we minstens een van beide waar maken.
We doen dan alsof een van beide in~$\Gamma$ staat.

Dit geeft ons een plan: voeg meer !!{formula}s aan~$\Gamma$ toe. Doe
dat zo dat (a)~de nieuwe verzameling consistent blijft; (b)~voor elke
mogelijke atomaire !!{sentence}~$!A$ ofwel $!A$ ofwel $\lnot !A$ in de
nieuwe verzameling staat; en (c)~bij $!A \land !B$ ook zowel $!A$ als
$!B$ in de verzameling staat, en bij $!A \lor !B$ ook minstens een van
$!A$ en $!B$, enzovoort. We blijven zulke formules toevoegen, zo nodig
eindeloos. Noem de verzameling van alle toegevoegde
!!{formula}s~$\Gamma^*$. De eerdere constructie geeft dan
\iftag{FOL}{!!a{structure}~$\Struct{M}$}{!!a{valuation}~$\pAssign{v}$}.
Met inductie kunnen we bewijzen dat zij alle zinnen in~$\Gamma^*$ waar
maakt. Zij maakt dan ook alle zinnen in~$\Gamma$ waar, want
$\Gamma \subseteq \Gamma^*$. Uiteindelijk blijken (a) en~(b) al
genoeg. Een verzameling zinnen die aan (b) voldoet, heet
\emph{volledig}. We moeten de consistente verzameling~$\Gamma$ dus
uitbreiden tot een verzameling~$\Gamma^*$ die zowel consistent als
volledig is.

\iftag{FOL}{%
Dit plan heeft nog een probleem. Stel dat
$\lexists[x][!A(x)] \in \Gamma$. Dan willen we een !!{constant}~$c$
kiezen en $!A(c)$ toevoegen. Maar is dat altijd mogelijk? Misschien
heeft onze taal maar enkele !!{constant}s en geldt voor elk daarvan
$\lnot !A(c) \in \Gamma$. We kunnen dan niet ook $!A(c)$ toevoegen:
de verzameling zou inconsistent worden. Bovendien weten we dan niet of
$\Struct{M}$ de formule $!A(c)$ of juist $\lnot !A(c)$ waar moet maken.
Het kan zelfs gebeuren dat $\Gamma$ alleen zinnen bevat in een taal
zonder constante symbolen, zoals de taal van de verzamelingenleer.

We lossen dit op door aan het begin oneindig veel constanten toe te
voegen. We voegen ook zinnen toe die deze constanten op de juiste
manier met de kwantoren verbinden. We moeten natuurlijk bewijzen dat
dit geen inconsistentie kan veroorzaken.

Onze eerste constructie werkt goed als de atomaire zinnen alleen
!!{constant}s bevatten. Maar een taal kan ook !!{function}s bevatten.
Dan moeten we geschikte functies op~$\Nat$ vinden om aan die
!!{function}s toe te kennen, en dat kan lastig zijn. We gebruiken
daarom een truc. Gebruik niet $i$ als betekenis van $\Obj c_i$, maar
neem de verzameling !!{constant}s zelf als domein. De structuur
$\Struct M$ kan dan elke !!{constant} aan zichzelf toekennen:
$\Assign{\Obj c_i}{M} = \Obj c_i$. We gaan nog een stap verder. Laat
$\Domain{M}$ bestaan uit alle \emph{termen} van de taal!{} Dan is
meteen duidelijk welke functie bij elk functiesymbool hoort. Deze
functies krijgen termen als invoer en leveren termen als uitvoer. Aan
de !!{function}~$\Obj f^n_i$ kennen we de functie toe die bij de $n$
termen $t_1$, \dots,~$t_n$ als invoer de term
$\Obj f^n_i(t_1, \dots, t_n)$ als uitvoer geeft.

Nu ontbreekt nog één puzzelstuk: wat doen we met~$\eq$? Het
!!{predicate}~$\eq$ heeft een vaste betekenis:
$\Sat{M}{\eq[t][t']}$ precies als
$\Value{t}{M} = \Value{t'}{M}$. In onze constructie is de !!{value} van
een term~$t$ de term $t$ zelf. Deze !!{structure} maakt daardoor
\emph{geen enkele} zin van de vorm $\eq[t][t']$ waar, behalve als $t$
en $t'$ letterlijk dezelfde term zijn. Dat is een probleem. Vrijwel
elke interessante theorie in een taal met !!{function}s heeft namelijk
zinnen $\eq[t][t']$ als stelling waarin $t$ en $t'$ verschillende
termen zijn. In een rekenkundige theorie is bijvoorbeeld
$\eq[(\Obj 0+ \Obj 0)][\Obj 0]$ een stelling. We passen daarom het
domein van~$\Struct M$ aan. De objecten in~$\Domain{M}$ zijn niet
langer losse termen, maar verzamelingen van termen. Elke verzameling
bevat precies de termen die volgens de zinnen in~$\Gamma$ gelijk moeten
zijn. Stel bijvoorbeeld dat $\Gamma$ een rekenkundige theorie is. Dan
bevat een van deze verzamelingen $\Obj 0$, $(\Obj 0 + \Obj 0)$,
$(\Obj 0 \times \Obj 0)$, enzovoort. We kennen die verzameling toe aan
$\Obj 0$. Zij blijkt ook de waarde te zijn van alle termen in de
verzameling, waaronder $(\Obj 0 + \Obj 0)$. In de aangepaste
!!{structure} is de zin $\eq[(\Obj 0+ \Obj 0)][\Obj 0]$ daarom waar.}{}

Dit wordt dus de volledige route. Eerst onderzoeken we consistente
verzamelingen die !!{complete} zijn. We bewijzen vooral het volgende:
zo'n !!a{complete} consistente verzameling bevat $!A \land !B$ precies
als zij zowel $!A$ als~$!B$ bevat; zij bevat $!A \lor !B$ precies als
zij minstens een van beide bevat; enzovoort
(\olref[ccs]{prop:ccs}).\iftag{FOL}{ Daarna definiëren en onderzoeken
  we ``verzadigde'' verzamelingen zinnen. Zo'n verzameling bevat
  implicaties die elke gekwantificeerde !!{sentence} koppelen aan haar
  concrete gevallen (\olref[hen]{defn:henkin-exp}). We bewijzen dat we
  elke consistente verzameling~$\Gamma$ kunnen uitbreiden tot een
  verzadigde verzameling~$\Gamma'$ (\olref[hen]{lem:henkin}). Is een
  verzameling consistent, verzadigd en !!{complete}? Dan geldt ook:
  zij bevat \iftag{prvEx}{$\lexists[x][!A(x)]$ precies als zij $!A(t)$
    bevat voor een gesloten term~$t$}{}\iftag{defEx,defAll}{}{ en
  }\iftag{prvAll}{$\lforall[x][!A(x)]$ precies als zij $!A(t)$ bevat
    voor elke gesloten term~$t$}{}
  (\olref[hen]{prop:saturated-instances}).}{}
Vervolgens nemen we de\iftag{FOL}{ verzadigde}{} consistente
verzameling~$\iftag{FOL}{\Gamma'}{\Gamma}$. We breiden haar uit tot een
\iftag{FOL}{verzadigde, consistente en}{consistente en}
!!{complete} verzameling~$\Gamma^*$ (\olref[lin]{lem:lindenbaum}). Met
$\Gamma^*$ definiëren we vervolgens ons \iftag{FOL}{termmodel
  $\Struct M(\Gamma^*)$}{!!{valuation}~$\pAssign v(\Gamma^*)$}.
\iftag{FOL}{Het domein van het termmodel bestaat uit de gesloten termen.
  De atomaire !!{sentence}s in de verzameling geven aan hoe we de
  !!{predicate}s moeten interpreteren}{De !!{propositional variable}s
  in de verzameling bepalen de valuatie} in~$\Gamma^*$
(\olref[mod]{defn:termmodel}). Daarna gebruiken we de eigenschappen van
\iftag{FOL}{ verzadigde,}{} volledige consistente verzamelingen. Zo
bewijzen we dat inderdaad
$\iftag{FOL}{\Sat{M(\Gamma^*)}}{\pSat{v(\Gamma^*)}}{!A}$ precies als
$!A \in \Gamma^*$ (\olref[mod]{lem:truth}). In het bijzonder volgt
dan $\iftag{FOL}{\Sat{M(\Gamma^*)}}{\pSat{v(\Gamma^*)}}{\Gamma}$.
\iftag{FOL}{Ten slotte bekijken we hoe we een termmodel definiëren als
  $\Gamma$ ook~$\eq$ bevat (\olref[ide]{defn:term-model-factor}). We
  bewijzen dat dit model~$\Gamma^*$ vervult (\olref[ide]{lem:truth}).}{}

\end{document}
